% ==============================================================================
% PHỤ LỤC A: HƯỚNG DẪN BIÊN DỊCH, KIỂM THỬ VÀ TÁI LẬP SỐ LIỆU
% ==============================================================================
\chapter{Hướng dẫn biên dịch, kiểm thử và tái lập số liệu}
\label{app:huong-dan}

\section{Biên dịch và chạy}

\begin{khungbash}[Các lệnh chính]
$ make                # biên dịch chương trình mẫu build/main
$ make run            # chạy main.cpp với dữ liệu mẫu của cả 5 task
$ make test           # kiểm tra thư viện được phép + 23 ca kiểm thử
$ make report-data    # sinh lại docs/report_data.txt (vết và bảng số của báo cáo)
$ scripts/package.sh <MSSV>   # tạo dist/<MSSV>.rar để nộp
\end{khungbash}

Biên dịch thủ công không cần \ma{make}:

\begin{khungbash}[Biên dịch trực tiếp]
$ g++ -std=c++11 -O2 -Isrc src/PathNode.cpp src/Algo.cpp src/main.cpp -o main
$ ./main
\end{khungbash}

\section{Tương thích với CodeRunner}

\begin{itemize}
  \item \ma{Algo.cpp} và \ma{PathNode.cpp} chỉ dùng \ma{<string>}, \ma{<vector>},
        \ma{<cmath>}, \ma{<algorithm>}; \ma{main.cpp} dùng thêm \ma{<iostream>}.
  \item Hàm \ma{printPath} chỉ được định nghĩa trong \ma{main.cpp}, không có trong mã nộp
        CodeRunner, đúng yêu cầu của đề.
  \item Các biến \ma{astar::expandedCount} và \ma{astar::eventLog} chỉ phục vụ đo đạc; khi
        \ma{eventLog} rỗng (mặc định), chi phí thêm là một phép so sánh con trỏ mỗi lần ghi.
\end{itemize}

\section{Đầu ra mẫu}

Trích đầu ra của \ma{make run} cho Task 4 (đầy đủ trong \ma{docs/sample_output.txt}):

\begin{khungbash}[make run (trích)]
=== Task 4: Game grid (0,0) -> (4,4) (Chebyshev) ===
Solution Path:
Node: Start | f: 4 | g: 0 | h: 4
Node: Right | f: 5 | g: 1 | h: 4
Node: Right | f: 6 | g: 2 | h: 4
Node: Right | f: 7 | g: 3 | h: 4
Node: Down-Right | f: 7.5 | g: 4.5 | h: 3
Node: Down | f: 7.5 | g: 5.5 | h: 2
Node: Down | f: 7.5 | g: 6.5 | h: 1
Node: Down | f: 7.5 | g: 7.5 | h: 0
\end{khungbash}
